\section{环同态}

记号约定:$A,B,C$是环.

\begin{definition}{环同态}{}
	设$f\in\Hom_{\mathsf{Set}}\left(A,B\right)$.若$f$是加法群的同态和乘法群的同态,则称$f$为环同态.记
	\begin{align*}
		\Hom_{\mathsf{Ring}}\left(A,B\right)=\left\{\varphi\in\Hom_{\mathsf{Set}}\left(A,B\right)\middle|\varphi\text{是环同态}\right\},\End_{\mathsf{Ring}}\left(A\right)=\Hom_{\mathsf{Ring}}\left(A,A\right).
	\end{align*}
	称$\End_{\mathsf{Ring}}\left(A\right)$的元素为$A$上的环同态.
\end{definition}

\begin{definition}{环同构}{}
	设$f\in\Hom_{\mathsf{Ring}}\left(A,B\right)$.若存在$g\in\Hom_{\mathsf{Ring}}\left(B,A\right)$使得$g\circ f=\id_{A}$且$f\circ g=\id_{B}$,则称$f$为环同构.记
	\begin{align*}
		\Isom_{\mathsf{Ring}}\left(A,B\right)=\left\{f\in\Hom_{\mathsf{Ring}}\left(A,B\right)\middle|f\text{是环同构}\right\},\Aut_{\mathsf{Ring}}\left(A\right)=\Isom_{\mathsf{Ring}}\left(A,A\right).
	\end{align*}
	称$\Aut_{\mathsf{Ring}}\left(A\right)$的元素为$A$上的环同构.
\end{definition}

\begin{proposition}{环同态的性质}{}
	设$f\in\Hom_{\mathsf{Ring}}\left(A,B\right),g\in\Hom_{\mathsf{Ring}}\left(B,C\right)$.
	\begin{enumerate}
		\item $g\circ f\in\Hom_{\mathsf{Ring}}\left(A,C\right)$.
		\item $f\in\Isom_{\mathsf{Ring}}\left(A,B\right)$ $\iff$ $f$是双射.两条件之一成立时$f^{-1}\in\Isom_{\mathsf{Ring}}\left(B,A\right)$.
		\item $f|_{A^{\times}}^{B^{\times}}$是乘法群之间的同态.
	\end{enumerate}
\end{proposition}

\begin{example}{}{}
	\begin{enumerate}
		\item $\Hom_{\mathsf{Ring}}\left(\Z,A\right)$中的元素只有$f:\Z\to A,n\mapsto n1_{A}$.特别地,$\Aut_{\mathsf{Ring}}\left(\Z\right)=\left\{\id_{\Z}\right\}$.
		\item 设$a\in A^{\times}$,则$\Ad_{a}:A\to A,x\mapsto axa^{-1}$是环同构,称为$a$诱导的内自同构,
	\end{enumerate}
\end{example}
































































































